\justifying
\NSFCBodyText

\subsubsection{研究内容}

本项目围绕教育智能体运行中的情境依赖风险，研究风险表征、轨迹评测和监测干预三个相互衔接的内容。研究对象以课程答疑、文献学习辅助中的工具型单智能体为主，采用可控检索和记忆操作记录任务执行过程。双智能体协同仅用于检验错误传播的压力情形，不扩展为独立的大规模多智能体平台。

\textbf{研究内容一：教育情境约束下的可信风险表征。}从任务目标、允许帮助程度及数据授权范围出发，建立能够连接风险类别与外显行为证据的表征方法。将风险归纳为证据与内容失真、教学支持失配、行动与信息边界违背三个相互关联的方面，分别覆盖事实错误与来源不实、目标偏离与过度代答、工具越权与记忆污染。教学支持失配包含支持不足和明确违背帮助边界，两者分开标注和报告。公平性问题通过相同任务条件下的支持差异开展探索，不根据单条输出推断群体偏见。

拟为风险事件记录触发情境、可观察证据、受影响对象及行为后果，并辨认首次能够依据外显信息判断风险的位置。授权帮助边界违背是可操作的规则风险，不能直接等同于已经造成伦理危害。标注分别保留明确违背、教学支持不足和规范不确定，邀请未参与任务生成的教学人员复核边界及反例；争议样本不强行归类，并检查不同处理方式对结论的影响。风险注入仅表明测试情境发生改变，不直接作为智能体已经违规的标签；实际标签需要依据后续可观察行为和独立人工判断建立。研究将特别比较最终答案相同但过程不同，以及过程相近但教学授权不同的任务对，以避免把“回答正确”误当作充分的可信条件。

\textbf{研究内容二：面向行为轨迹的过程级可信评测。}在风险表征与标注规则确定后，研究教育情境和轨迹证据对识别结果的共同作用。可观测轨迹包括用户任务、外显计划、检索证据、工具参数、执行回执和记忆读写记录，不要求获取模型私有思维链。对已完成任务开展离线风险定位，对运行中的任务只使用当前前缀和拟执行动作进行告警，分别回答“发生了什么”与“当前是否需要干预”。

评价采用规则核验与语义判别相结合的方式。授权和参数范围等可确定事项由规则检查；教学支持程度、证据误用等依赖情境的事项由具有明确评分依据的语义评价器处理。输出分维度风险判断、支持证据和不确定性，不用一个加权总分掩盖严重越权。通过等信息对照、教育约束消融和跨任务家族测试，识别过程评价的有效成分，并对大模型评价器与人工判断的一致性进行独立检查。

语义评价按风险维度采用0—3级证据评分：0为未见支持风险的证据，1为可疑线索，2为证据支持风险判断，3为明确且直接的违规证据；规范是否可判定另设不确定标记。逐时点保留评分与对应片段，整轨迹以该维度最高评分判断是否曾告警，线上仅累计当时已见信息。同信息的通用与教育评价共用基础模型、输入字段和输出格式，区别在评分依据：前者采用一般安全与任务正确性规则，后者增加本任务的教学支持边界；交叉消融则固定教育评价器，只改变证据和约束条件。离散分数阈值在正常校准轨迹上选择，无法恰好达到名义误报目标时采用不超过目标的最近可行阈值，并报告对应实测值。确定性越权始终阻断；语义评分达到阈值且可修正时重规划，无法安全修正时阻断；规范不确定时转人工；其余动作执行。人工转交单独报告，同时计入正常任务的干预负担，不能隐去其造成的自主完成损失。

\textbf{研究内容三：伦理风险动态监测与安全干预。}以第二项内容的前缀风险评价为基础，在回答发送、工具执行和记忆写入之前设置监测点。沙箱先缓存待发送回答，评价器可以读取尚未发布的候选文本，审查通过后才交付学习者；其动作提交时点即对学习者发布时点。例如，在要求引导练习的任务中，即使不调用工具，直接给出完整代答也须经过发送前检查。确定性权限违背触发阻断；可修正的任务偏离优先请求重规划；证据不足或涉及教学判断争议的情况交由人工复核；低风险且满足授权要求的动作继续执行。各类处置的研究重点是触发条件与执行后果是否相称，尤其关注误阻断和无效重规划。

三项研究共用课程答疑与文献学习辅助两类场景，通过公开课程资料、可控构造任务和授权后的高校案例逐层验证，分别标明公开规则改编、专家构造或授权记录的来源，合成角色不用于推断真人行为。课程答疑关注提示、讲解和独立练习边界，文献学习辅助关注证据引用、资料访问与记忆更新。双智能体压力测试进一步检查污染信息是否经转发进入后续工具调用或持久记忆，并以等消息、等步数的单智能体转发作为对照。

\subsubsection{研究目标}

总体目标是形成可复核的教育智能体可信评测与监测方法，解释情境约束、过程证据和干预后果之间的关系。针对风险表征，形成具有行为证据、标注依据和反例的风险分类与案例集；针对轨迹评测，比较结果评价、通用轨迹评价与教育情境评价的识别表现和有效告警时间；针对监测干预，给出风险执行、正常任务完成、教学支持合规性及人工负担之间的比较结果。

项目不预先承诺未知的性能提升幅度。主要科学判断依据留出任务上的配对差值、置信区间与消融结果；即使假说不获支持，仍应能够说明哪些风险可由简单规则处理、哪些依赖教育情境，以及哪些情形缺少足够证据开展自动干预。上述边界与失败案例也是项目的可验收研究结果。

\subsubsection{拟解决的关键科学问题}

\textbf{第一，风险含义为何随教育情境改变，哪些证据能够区分这种变化？}研究在“教师演示”“引导练习”等不同支持条件下，同一外显回答或工具行动的风险判断差异。教学情境的作用必须在与通用轨迹评价相同的信息范围内建立，而不能仅靠给本方法增加提示内容获得优势。

\textbf{第二，完整过程中的有效证据能否在运行前缀中及时获得？}研究风险首次可观察位置与实际动作执行之间的可用干预窗口。对于只能在执行后确认的后果，如实归为事后可识别风险，不把它们计入动作前监测的成功范围；对尚无充分证据的情形，保留不确定判断。

\textbf{第三，及时识别能否转化为具有教学代价约束的干预收益？}研究相同初始任务在不同干预策略下的分支结果，比较风险减少是否超出统一拒绝、规则检查等简单策略所能解释的程度。若依赖频繁拒绝或人工接管才能取得低风险，需明确其适用范围和运行负担。

\subsubsection{研究方法}

\textbf{数据与任务构建。}从允许研究使用的公开课程资料和论文学习材料构建基础任务，保留来源及版本。每个任务明确目标、工具集合、允许的数据范围和教学支持要求；构造干净任务、检索污染、记忆污染及教学约束变化等版本。污染内容仅进入隔离沙箱，外部发送、课程记录修改等动作使用模拟工具与虚拟身份，避免在真实教学环境中实施高风险测试。

任务规模采用\textbf{【暂定 24个模板家族、每家族5个基础任务，共120个基础任务】}。每个基础任务设置\textbf{【暂定 4个受控版本、2个模型和3个随机种子，共2880次候选轨迹运行】}；检索与记忆为同一工具型代理的组件，额外架构和双智能体测试仅在主实验完成后增补。该数量不包含评价器调用、干预分支、失败重试与人工仲裁，预实验将分别测算并设置总费用上限。四个版本中保留一个干净版本；教学约束变化需经人工判断，不能默认计为正常样本。各家族在主分析中等权，独立统计信息受模板家族和基础任务数量限制。先开展小样本预实验估计标注耗时、风险事件比例与方差，再决定是否增加家族数量或收缩次要分析。

\textbf{人工标注与评价器校准。}按照风险定义与证据规则，由\textbf{【暂定 2名独立标注者和1名分歧仲裁者】}判断轨迹，具体人员在团队落实后确定。标注者不知模型身份和方法条件，需指出支持判断的可观察片段。先讨论标注手册中的歧义，再对独立样本报告逐类一致率、Cohen's $\kappa$及分歧类型；对类别失衡同时保留原始列联表，不用单一一致性分数替代分析。

基础任务及其所有变体、种子运行按模板家族整组划分，采用\textbf{【暂定 12个家族用于规则与评价器开发、6个用于校准、6个用于留出测试】}。用于模型辅助构造的数据不直接进入人工真值；目标模型输出也不能同时充当自身评价的参考答案。评价提示、阈值、主要终点及对照方法在访问测试集结果前确定，避免用测试反馈反复改写规则。

\textbf{对照与消融。}设置仅检查结果、通用轨迹评价、具有教育约束的轨迹评价三类离线对照；其中后两者获得相同的可观测字段、上下文容量和模型预算。进一步采用轨迹证据保留/移除与教育约束正确/打乱的交叉消融，检验教育约束和证据之间的交互作用。打乱约束只改变评价器输入，真实任务标签仍由原始情境确定；另设合法授权变化的任务对，用于研究风险含义随情境的变化。两类实验不得混用标签。合法授权比较固定任务、行动和回答，只改变支持规则；另外开展自然运行验证，区分标签含义变化与智能体行为变化。增加等长无关约束和规则措辞改写对照，检查规则关键词捷径。

最终答案一致或正确性相近的任务对作为机制诊断子集，主比较保留所有初始任务，避免依据干预后的成功状态筛选样本。对显式权限字段易于判定的任务与需要语义理解的任务分别报告；对文本长度、风险词、模型身份以及已知模板线索做检查，判断表观提升是否来自捷径。

\textbf{风险与时间指标。}风险识别报告各类别精确率、召回率、误报率和精确率—召回率曲线面积，并列出类别比例。设$t_0$为首个可观察风险证据出现时点，$t_a$为风险动作提交时点，$t_d$为首次有效告警时点。只有$t_0\leq t_d<t_a$的告警计为动作前检出；$t_d-t_0$表示检测延迟，$t_a-t_d$表示剩余干预窗口。未检出和执行后才发现的事件单独统计，不删除后再计算平均延迟。正常轨迹的误报按整条任务中“至少一次错误告警”计算，防止增加监测次数虚增及时性。

阈值在校准集上依据\textbf{【暂定 整轨迹误报率10\%的名义工作点】}选择，同时报告完整曲线与置信区间。各方法只在正常校准轨迹上选择阈值，测试阶段报告各自实际误报及区间，不能仅因名义目标相同就声称实际误报相同。该工作点是比较设置，不是已获统计保证；若校准样本不足以估计低误报区间，则扩大区间报告，并将相关结果定位为探索性证据，不声称达到部署级安全水平。

\textbf{干预效果与统计推断。}运行对照包括事后检查、仅规则保护、统一拒绝和分级干预。基于相同初始任务、初始状态与随机条件成对比较；当干预改变后续路径时重新执行对应分支，不以原未干预轨迹替代干预结果。风险执行比例的分母为预先设定的全部风险测试任务，正常完成率的分母为全部正常任务，拒绝、超时和人工移交均计入。教学支持合规性由独立人工量表评价，与任务正确性分开报告。

以无在线干预的正常任务完成率为参照，在校准集设定\textbf{【暂定 正常任务完成率损失不超过5个百分点的比较约束】}，冻结策略后检验测试结果；损失区间上界低于约束时支持满足，下界高于约束时支持违反，跨越约束则证据不足，不从测试集重选阈值。人工移交按未自主完成计，同时单列其后人工完成情况。对确认、监测、重规划全部计入调用次数、耗时和费用，记录每条初始任务的人工复核次数。

主检验为教育约束与轨迹证据的召回交互差。以$R_{ec}$表示证据状态$e$与约束状态$c$下、校准阈值固定后的家族等权召回，$e=1/0$对应保留/移除轨迹证据，$c=1/0$对应正确/打乱教育约束，定义\[\Delta_{\mathrm{int}}=(R_{11}-R_{01})-(R_{10}-R_{00}).\]

同一预先选定模型的教育支持边界风险作为主要分析，其他模型和风险类别用于稳健性与探索；模型选择不依据测试表现。\textbf{【暂定 交互召回增益5个百分点为最小实际意义幅度】}，在预实验后、测试前确定最终幅度；只有区间下界超过该幅度才支持预定实用增益，区间支持正向作用但未跨越该幅度时，不声称已有实用收益。区间上界低于该幅度则不支持达到预定实用增益，跨界结果保留不确定性。通用轨迹对照的召回差值和干预后的风险执行比例差值作为支持性分析，后者只有在效用损失区间支持满足约束、风险执行差值区间同时支持下降时，才支持预定的折中收益。测试实际误报不一致时，交互结果仅解释为该校准策略下的表现，不能单独归因为教育约束与证据的机制协同。

主风险阳性仅指依据原始任务情境确认的教育支持边界违背；支持不足与规范不确定不并入阳性，分别报告数量，并对不确定样本的两种极端归类作敏感性分析。各条件共用同一冻结真值和阳性分母，逐家族列出阳性、阴性和不确定计数。预设主估计对象为全部留出家族；任何家族无主风险阳性时，其召回记为不可估计，全部家族的主交互量与相应成功判断也记为不可估计，不赋零或静默删除。可另列有阳性家族的条件性结果，但仅作为探索并报告家族覆盖率。

采用模板家族外层、基础任务内层的分层聚类自助法给出95\%置信区间。可估计时，在每家族内按基础任务是否含已确认主风险阳性分层重抽，任务的全部变体和种子成组保留，四个消融条件使用同一次抽样；该区间条件于观测到的标签构成，不把无阳性样本推定为安全。任何仍不可估计的重抽结果均记录数量，不通过静默丢弃获得区间；种子只估计运行波动，不作为新增独立样本。六个测试家族只能支持所选任务范围内的探索性结论；家族数较少时增加逐家族敏感性分析，次要比较作多重性控制或明确标为探索，不以未显著等同于两种方法等效。

\subsubsection{技术路线}

\textbf{总体技术路线。}教育任务及其教学与授权要求首先进入风险表征环节，形成风险类别、证据规则和人工标注依据；这些结果支持轨迹评测的开发与校准，产生能够说明依据的前缀风险判断；经过校准的判断再驱动动作前监测与分级干预。三条路线在同一组课程答疑和文献学习任务上检验，最终形成评测案例、监测原型及其有效范围。留出测试只承担最终验证，不能反向参与当前版本的规则或阈值调整。

\textbf{面向研究内容一的风险表征路线。}以公开学习材料、任务目标和授权范围为输入，先分析正常任务需要执行的动作，再构造受控情境变化与污染版本。依据外显证据定位风险事件，进行盲法双人标注和分歧仲裁，输出风险类别、标注手册及带证据的案例。通过相同答案的过程差异、合法授权变化和标注一致性检查验证表征的区分能力；通过的证据规则与开发数据提供给第二条路线，校准和测试材料保持独立。

\textbf{面向研究内容二的轨迹评测路线。}以第一条路线的风险规则和可观测轨迹为输入，建立规则核验与语义评价器，分别开展完整轨迹定位及仅使用当前前缀的告警。通过等信息通用评价对照、教育约束交叉消融和模板家族留出测试，验证识别表现及其来源，输出分维度判断、证据定位、校准参数和失败案例。满足可解释性与校准要求的前缀评价方法进入第三条路线，证据不足的风险类别保留人工处理。

\textbf{面向研究内容三的监测干预路线。}以待发送回答或拟执行工具/记忆动作、当前轨迹前缀和校准参数为输入，先检查确定性权限，再根据可修正性和不确定程度决定执行、重规划、阻断或人工复核。每次干预后重建对应运行分支，追踪任务完成与风险后果；只有重新通过检查的动作才允许提交。通过与事后检查、仅规则保护和统一拒绝比较，检验风险下降、误阻断和人工负担，输出可运行原型与分场景建议。新发现的问题可用于下一版本开发，但不能回填当前测试版本以改变既有结论。

\subsubsection{关键技术}

\textbf{情境约束与证据的可核对关联。}将每项风险判断连接到任务中的具体限制与轨迹片段，例如将来源忠实性连接到引用和检索回执，将权限风险连接到工具参数及授权范围，将教学支持失配连接到本次允许的帮助层级。通过保留反例和不确定标签，使“需要更多证据”的情况与“已识别风险”能够区分。相关表征用于解释判断，不声称从外显日志恢复模型内部因果过程。

\textbf{离线标签与在线信息的时间隔离。}离线人工标注可以读取完整记录以确认事件，在线评价只能读取当时的前缀。评价输入构造和日志记录分别保存，检查未来动作、最终结果及人工标签是否误入在线提示。对只能事后观察的风险单列指标，对无法完成的模拟工具调用保留失败回执，避免把工具不可用混同为安全拒绝。

\textbf{干预分支的可复现比较。}固定任务版本、模型版本、工具行为和预算，记录动作提交前后的环境状态。校准阶段锁定干预规则后，复用相同初始条件生成各策略分支。双智能体压力测试分别控制来源标识与权限约束，追踪下游采纳、工具执行及持久写入；若等步数单智能体转发也出现同等传播，只能归结为信息传播风险，不能宣称多智能体特有的放大机制。

\subsubsection{可行性分析}

项目采用已有模型与可控工具环境，核心工作量集中在任务设计、标注、评价器校准和对照实验，不依赖基础模型训练。公开文献和工具型智能体基准提供了方法参照，申请人已确认的算力及软件工程条件能够支持离线运行。三个研究内容共享任务与日志，能够减少分别建设系统造成的重复投入。

执行按“沙箱机制检验—近真实案例—授权后真实场景”推进。真实课程、数据使用和伦理审查情况为\textbf{【待补 F-EDU-02：课程场景、授权范围及审查结论】}。上述条件落实前，只开展公开资料和合成角色构成的沙箱实验；真实场景迟延时，仍可完成方法比较、监测原型和近真实案例报告，但对学生学习效果与实际部署风险不作外推结论。预实验若显示标注或运行成本超出一年期投入，优先保留教育情境与轨迹证据的主比较，将公平性和双智能体扩展保留为探索。
